\documentclass[a4paper,12pt]{report}
\pdfminorversion=4
\usepackage[utf8]{inputenc}
\PassOptionsToPackage{hidelinks}{hyperref}
\usepackage[a-1b]{pdfx}

%  PREAMBOLO: stesso impianto della tesi (tesi.sty, frontespizio ridefinito)
\usepackage[utf8]{inputenc}
\usepackage[italian]{babel}
\usepackage{setspace}
\usepackage{graphicx}
\usepackage{tesi}
\usepackage[a4paper]{geometry}

%  FRONTESPIZIO: identico a quello di main.tex
\makeatletter
\def\@sottotitolo{}
\newcommand{\sottotitolo}[1]{\gdef\@sottotitolo{#1}}
\def\titlep{%
  \thispagestyle{empty}%
  \begin{center}
    \includegraphics[width=0.78\textwidth]{img/logo_unimi.png}\\[0.8em]
    \Large {\em \@dept}
  \end{center}
  \vskip0.5in
  \begin{center}
    {\hyphenpenalty=10000\setstretch{1.15}\huge\bfseries \@title\par}%
    \ifx\@sottotitolo\@empty\else
      \vspace{1.4em}%
      {\hyphenpenalty=10000\setstretch{1.15}\fontsize{15}{18}\selectfont \@sottotitolo\par}%
    \fi
    \vspace{2.5em}%
    {\Large\sc Riassunto della tesi\par}%
  \end{center}
  \vfill
  \begin{description}
    \item[Relatore:] \@relatore
    \item[Correlatore:] \@correlatore
  \end{description}
  \vskip0.5in
  \null\hfill
  \parbox{3in}{Tesi di: \\ \@author \\ Matricola: \@matricola}
  \vfill
  \begin{center}\large Anno Accademico \@anno\end{center}
  \newpage}
\makeatother

\begin{document}
\title{Vulnerabilità avversariali nei sistemi d'arma autonomi tra algoritmi e diritti}
\sottotitolo{Un'analisi sperimentale sulla robustezza della computer vision, implicazioni per il principio di distinzione e per il controllo umano}
\author{Daniele Intra}
\dept{Corso di Laurea Magistrale in Sicurezza Informatica}
\anno{2025-2026}
\matricola{53665A}
\relatore{Prof. Matteo Buffa}
\correlatore{Prof. Massimo Walter Rivolta}

\titlep
%  Testo: indicazioni di Ateneo, margini 2.5 cm e interlinea 1.5
\newgeometry{margin=2.5cm}
\pagestyle{plain}
\setcounter{page}{1}
\onehalfspacing

I sistemi d'arma autonomi letali (Lethal Autonomous Weapon Systems, LAWS)
affidano a un algoritmo la selezione del bersaglio, un compito che finora
spettava a un operatore umano. Il Diritto Internazionale Umanitario (DIU) impone
di distinguere fra combattenti e civili e di mantenere un controllo umano
significativo sull'uso della forza. Questi principi presuppongono però che il
sistema percepisca correttamente la scena. Nei LAWS la percezione è affidata a
reti neurali di visione artificiale, note per essere vulnerabili agli attacchi
avversariali: alterazioni dell'immagine costruite appositamente possono far
sparire una persona agli occhi del rilevatore. Il problema affrontato dalla tesi
è quanto sia affidabile questo canale percettivo e che cosa accada alla
decisione di ingaggio quando viene manipolato.

La questione non è teorica. Sistemi con autonomia letale sono già stati
impiegati, per esempio con il drone Kargu-2 in Libia e nella guerra in Ucraina,
e piattaforme di targeting come Lavender e Habsora riducono la distinzione fra
combattente e civile alla scelta di una soglia su un punteggio. La prima parte
della tesi mostra che le regole non tengono il passo: i negoziati nell'ambito
della Convenzione su certe armi convenzionali (CCW) sono fermi, l'AI Act europeo
esclude gli usi militari e le verifiche statunitensi restano interne alla catena
di comando. Anche i vincoli etici fissati dalle aziende cedono alle esigenze
degli Stati, come mostra lo scontro fra Anthropic e il Pentagono. In questo
quadro nessuno misura in modo indipendente l'affidabilità del sensore su cui
poggia il principio di distinzione.

La ricerca sugli attacchi avversariali ha studiato soprattutto persone riprese
da terra, che occupano gran parte dell'immagine. Nei pochi lavori su immagini
aeree i risultati più solidi riguardano i veicoli, e in alcuni casi i pedoni
sono esclusi perché troppo piccoli. Manca quindi una misura, con un protocollo
statistico esplicito, proprio per la classe che il DIU impone di proteggere: la
persona ripresa da un drone.

Per colmare questa lacuna è stato sviluppato LAWS-SIM, un framework articolato
in tre livelli. Nel primo si ottimizza, contro il rilevatore di persone YOLOv8
e con piena conoscenza della rete (regime \emph{white-box}), una patch
avversariale universale: una trama che copre il $20\%$ della figura sulla parte
superiore del tronco e può quindi essere stampata su un indumento. Perché
resti efficace in condizioni di ripresa variabili, la patch è ottimizzata con
la tecnica Expectation over Transformation (EoT), cioè in media su sedici
versioni ruotate, ridimensionate e diversamente illuminate per ogni
fotogramma, e con un vincolo di variazione totale che ne limita le alte
frequenze e la rende stampabile. L'effetto è misurato su due insiemi di dati
aerei indipendenti, VisDrone2019-DET e Okutama-Action, confrontando gli stessi
fotogrammi con e senza patch. La metrica principale è l'Evasion Rate, la
frazione di fotogrammi contenenti una persona in cui il rilevatore non ne
individua alcuna; la significatività statistica dell'uso della patch è stabilita con un
bootstrap appaiato su
$10\,000$ ricampionamenti, dopo aver mitigato la correlazione fra fotogrammi
consecutivi delle sequenze video. Il secondo e il terzo livello formano un
simulatore descritto come macchina a stati, in cui il canale visivo,
parametrizzato con le misure del primo livello, è combinato con informazioni da
fonti aperte (Open Source Intelligence, OSINT) e con l'analisi del
comportamento, fino alla decisione di ingaggio.

La patch aumenta l'Evasion Rate in modo statisticamente significativo e con
entità quasi identica sui due domini: dal $17.5\%$ al $42.5\%$ su VisDrone
($+25.0$ punti percentuali) e dal $49.7\%$ al $77.7\%$ su Okutama-Action
($+28.0$ punti), con $p < 0.0001$ in entrambi i casi. L'effetto riguarda solo
chi indossa la patch: le altre persone nella scena restano rilevate e i falsi
allarmi non aumentano. Sei configurazioni di ottimizzazione molto diverse
raggiungono lo stesso limite, che dipende dalle dimensioni ridotte della figura
umana vista dall'alto più che dall'ottimizzatore. La stessa fragilità esiste
anche senza attacco: modificando la sola soglia di confidenza del rilevatore,
l'Evasion Rate su VisDrone passa da meno del $4\%$ a oltre il $72\%$. Nel
simulatore la patch da sola è compensata dagli altri canali; quando anche le
fonti aperte sono manipolate, invece, la quota di bersagli segnalati scende dal
$34\%$ al $15\%$, nella stessa proporzione misurata sul rilevatore, e il $96\%$
dei civili viene segnalato come sospetto.

Ne emerge che un sistema chiamato a riconoscere persone da quota lavora con un
margine ridotto già in condizioni normali, e che una contromisura pensata per
essere stampata su un capo di abbigliamento basta a modificarne l'esito.
L'operatore che deve approvare le decisioni del sistema non ha modo di
accorgersi della manipolazione, e il controllo umano significativo rischia di
restare una formalità. La sicurezza informatica del canale percettivo diventa
così una condizione tecnica del principio di distinzione. I risultati hanno
limiti dichiarati: l'attacco è stato valutato su immagini digitali e in regime
\emph{white-box}, per cui le misure sono un limite superiore dell'efficacia
ottenibile, e il simulatore ha valore illustrativo. I passi successivi sono la
prova fisica della patch stampata, la verifica su altri rilevatori e lo studio
di un dispositivo di autotutela per i civili, il C.A.R.E. Kit, da valutare
anche sul piano etico. Il codice è disponibile all'indirizzo
\url{https://github.com/intradaniele/Tesi-Laws}.

\end{document}
